\documentclass[10pt]{article}
\usepackage{amsmath,amssymb}
\usepackage[margin=1in]{geometry}
\begin{document}

\section*{Pinning-field observables and crossing field}

We use the signed field convention $h>0$ for plaquette pinning and $h<0$
for dimer--plaquette pinning.  For each rank $r$, repeated runs at fixed
$(J_2,D,h)$ are first averaged,
\begin{equation}
 \bar C_{r,D}(J_2,h)=\frac{1}{n_D}\sum_{i=1}^{n_D}C^{(i)}_{r,D}(J_2,h),
\end{equation}
so every retained bond dimension has the same total weight, irrespective of
the number of replicas or protocols available at that $D$.

For the three sorted antiferromagnetic nearest-neighbour correlations,
\begin{equation}
 C_{\rm strong}\le C_{\rm mid}\le C_{\rm weak},
\end{equation}
we define
\begin{equation}
 \Delta_D=C_{\rm weak}-C_{\rm strong},\qquad
 f_D=\frac{C_{\rm mid}-C_{\rm strong}}
           {C_{\rm weak}-C_{\rm strong}},\qquad
 q_D=1-2f_D
 =\frac{C_{\rm weak}+C_{\rm strong}-2C_{\rm mid}}
        {C_{\rm weak}-C_{\rm strong}}.
\end{equation}
Thus $q=-1$ is dimer--plaquette, $q=+1$ is plaquette, and $\Delta\ge0$
measures the total bond splitting.

At each sampled $(J_2,h)$, the three thermodynamic correlations are obtained
independently from the unweighted two-parameter fits
\begin{equation}
 \bar C_{r,D}(J_2,h)
 =C_{r,\infty}(J_2,h)+k_r(J_2,h)e^{-a_g(J_2)D},
 \qquad r\in\{\mathrm{strong},\mathrm{mid},\mathrm{weak}\},
\end{equation}
where $a_g(J_2)$ is fixed by the original-2C3 gapped energy fit.  Every
contiguous high-$D$ window containing the two largest available bond
dimensions is tested, using the same window for all three ranks.  We select
the window that minimises
\begin{equation}
 R=\max_r\left|C_{r,\infty}-\bar C_{r,D_{\max}}\right|,
\end{equation}
thereby rejecting formally exact but physically overshooting two-point
extrapolations.  Fit residuals are retained only as diagnostics.  Only after
these fits do we form
\begin{equation}
 \Delta(J_2,h)=C_{\rm weak,\infty}-C_{\rm strong,\infty},\qquad
 q(J_2,h)=
 \frac{C_{\rm weak,\infty}+C_{\rm strong,\infty}
       -2C_{\rm mid,\infty}}
      {C_{\rm weak,\infty}-C_{\rm strong,\infty}}.
\end{equation}
At $h=0$, the selected dimer--plaquette and plaquette data are fitted
separately by the same procedure, giving six extrapolated correlations.  The
corresponding ranks are then averaged,
\begin{equation}
 C_{r,\infty}(J_2,0)=\frac{C^{(d)}_{r,\infty}(J_2,0)
                              +C^{(p)}_{r,\infty}(J_2,0)}{2},
\end{equation}
and the preceding definitions give $\Delta(J_2,0)$ and $q(J_2,0)$.  No value
of either observable is imposed at $h=0$.

For the energy, the median-reduced data of each sector $s\in\{p,d\}$ and
each bond dimension are fitted by ordinary least squares,
\begin{equation}
 \widehat E_{s,D}(J_2,h)
 =a_{s,D}(J_2)+b_{s,D}(J_2)h+c_{s,D}(J_2)h^2.
\end{equation}
The final equal-$D$ constant extrapolation is
\begin{equation}
 E_s(J_2,h)=\frac{1}{N_D}\sum_{D\ge6}\widehat E_{s,D}(J_2,h),
 \qquad
 E(J_2,h)=
 \begin{cases}
 E_d(J_2,h),&h<0,\\
 E_p(J_2,h),&h>0.
 \end{cases}
\end{equation}
The crossing field is the real root closest to the origin,
\begin{equation}
 E_p(J_2,h_c)=E_d(J_2,h_c),\qquad |h_c|\le10^{-2}.
\end{equation}
For its uncertainty, the same equation is solved separately at every $D$,
giving $h_{c,D}$.  With
\begin{equation}
 s_h^2=\frac{1}{N_D-1}\sum_D
 \left(h_{c,D}-\overline h_c\right)^2,
\end{equation}
the plotted horizontal interval is
\begin{equation}
 h_c\ \pm\ 1.96\,\frac{s_h}{\sqrt{N_D}}.
\end{equation}

\end{document}
